% from ega1-6-fr.tex
\oldpage[I]{144}
\begin{proposition}[6.2.2]
\label{I.6.2.2-fr}
Étant donné un préschéma $X$, les conditions suivantes sont équivalentes~:
\begin{enumerate}
  \item[a)] $X$ est un schéma artinien~;
  \item[b)] $X$ est noethérien et son espace sous-jacent est discret~;
  \item[c)] $X$ est noethérien et les points de son espace sous-jacent sont
    fermés (condition $\mathrm{T}_1$).
\end{enumerate}

Lorsqu'il en est ainsi, l'espace sous-jacent à $X$ est fini, et l'anneau $A$
de $X$ est composé direct des anneaux locaux (artiniens) des points de $X$.
\end{proposition}

On sait que a) entraîne la dernière assertion (\cite{I-13}, p.~205, th.~3),
tout idéal premier de $A$ est alors maximal et est l'image réciproque de
l'idéal maximal de l'un des composants locaux de $A$, donc l'espace $X$ est
fini et discret~; a) entraîne donc b) et b) entraîne évidemment c). Pour voir
que c) entraîne a), montrons d'abord que $X$ est alors fini~; on peut en effet
se ramener au cas où $X$ est affine, et on sait qu'un anneau noethérien dont
tous les idéaux premiers sont maximaux est artinien (\cite{I-13}, p.~203),
d'où notre assertion. L'espace sous-jacent $X$ est alors discret, somme
topologique d'un nombre fini de points $x_i$, et les anneaux locaux
$\mathscr{O}_{x_i}=A_i$ sont artiniens~; il est clair que $X$ est isomorphe au
schéma affine spectre premier de l'anneau $A$ composé direct des $A_i$
(\hyperref[I.1.7.3-fr]{1.7.3}).

\subsection{Morphismes de type fini.}
\label{subsection:I.6.3-fr}

\begin{definition}[6.3.1]
\label{I.6.3.1-fr}
On dit qu'un morphisme $f:X\to Y$ est \emph{de type fini} si $Y$ est réunion
d'une famille $(V_\alpha)$ d'ouverts affines ayant la propriété suivante~:
\begin{enumerate}
  \item[(P)] $f^{-1}(V_\alpha)$ est réunion finie d'ouverts affines
    $U_{\alpha i}$ tels que chacun des anneaux $A(U_{\alpha i})$ soit une
    algèbre de type fini sur $A(V_\alpha)$.
\end{enumerate}
On dit alors aussi que $X$ est un préschéma de type fini sur $Y$, ou un
$Y$-préschéma de type fini.
\end{definition}

\begin{proposition}[6.3.2]
\label{I.6.3.2-fr}
Si $f:X\to Y$ est un morphisme de type fini, tout ouvert affine $W$ de $Y$
possède la propriété (P) de la déf. (\hyperref[I.6.3.1-fr]{6.3.1}).
\end{proposition}

Nous démontrerons d'abord le

\begin{lemma}[6.3.2.1]
\label{I.6.3.2.1-fr}
Si $T\subset Y$ est un ouvert affine, possédant la propriété (P), alors, pour
tout $g\in A(T)$, $D(g)$ possède la propriété (P).
\end{lemma}

En effet, par hypothèse, $f^{-1}(T)$ est réunion finie d'ouverts affines
$Z_j$ tels que $A(Z_j)$ soit une algèbre de type fini sur $A(T)$~; soit
$\varphi_j:A(T)\to A(Z_j)$ l'homomorphisme d'anneaux correspondant à la
restriction de $f$ à $Z_j$ (\hyperref[I.2.2.4-fr]{2.2.4}), et posons
$g_j=\varphi_j(g)$~; on a alors
$f^{-1}(D(g))\cap Z_j=D(g_j)$
(\hyperref[I.1.2.2.2-fr]{1.2.2.2}). Or,
$A(D(g_j))=A(Z_j)_{g_j}=A(Z_j)[1/g_j]$ est de type fini sur $A(Z_j)$ et
\emph{a fortiori} sur $A(T)$ en vertu de l'hypothèse, donc aussi sur
$A(D(g))=A(T)[1/g]$, ce qui démontre le lemme.

Ce lemme étant établi, comme $W$ est quasi-compact
(\hyperref[I.1.1.10-fr]{1.1.10}), il existe un recouvrement fini de $W$ par
des ensembles de la forme $D(g_i)$, où chaque $g_i$ appartient à un anneau
$A(V_{\alpha(i)})$. Chaque $D(g_i)$, étant quasi-compact, est réunion finie
d'ensembles $D(h_{ik})$ où $h_{ik}\in A(W)$~; si
$\varphi_i:A(W)\to A(D(g_i))$ est l'application canonique, on a
$D(h_{ik})=D(\varphi_i(h_{ik}))$ en vertu de
(\hyperref[I.1.2.2.2-fr]{1.2.2.2}). En vertu de
(\hyperref[I.6.3.2.1-fr]{6.3.2.1}), chacun des
$f^{-1}(D(h_{ik}))$ admet un recouvrement fini par des ouverts affines
$U_{ijk}$ tels que $A(U_{ijk})$ soit une algèbre de type fini sur
$A(D(h_{ik}))=A(W)[1/h_{ik}]$, d'où la proposition.

